% ---------------------------------------------------------------------------
%  Section 1 -- Introduction   (~540 words, one LNCS page)
%
%  Flow: problem -> what current methods achieve -> what they leave open ->
%  our observation -> contributions.
%
%  No numerical claims appear here, so this section does not need rewriting
%  when the final results land.
% ---------------------------------------------------------------------------

\section{Introduction}
\label{sec:intro}

Segmenting gliomas on magnetic resonance imaging (MRI) supports diagnosis,
surgical planning and the assessment of treatment response. The standard
protocol acquires four complementary sequences: T1, contrast-enhanced T1
(T1CE), T2 and FLAIR. No single sequence reveals the whole lesion, since
enhancement marks the active core while oedema is clearest on FLAIR. Manual
delineation is slow and varies between readers, which has made automated
segmentation a central problem in medical image analysis~\cite{brats_menze}.

Deep encoder--decoder architectures~\cite{unet,unet++,nnunet} perform well on
public benchmarks. A parallel line of work addresses an obstacle these models
ignore: the complete protocol is often unavailable, because examinations are
shortened, sequences are discarded for motion artefacts, or contrast is
contraindicated. Methods for incomplete multimodal segmentation respond by
learning representations defined for any available subset~\cite{hemis,rfnet},
by distilling from a full-modality teacher~\cite{mmformer}, or by separating
representation learning from fusion~\cite{unime}.

These methods share an assumption that is rarely stated: that the output
should be a single deterministic mask. Their evaluations report how accuracy
degrades as sequences are removed, which shows that the prediction depends on
the available evidence. The prediction itself, however, is returned in the
same form either way. A boundary inferred from FLAIR alone looks no different
from one inferred from all four sequences, so the clinician has no indication
of how confident the model is.

Generative methods model a distribution over plausible masks rather than one
answer~\cite{probunet,phiseg,medsegdiff}, and the spread across samples shows
what the model finds ambiguous. Validating that spread is harder than
producing it, because the quantity it estimates is rarely observable. The
usual solution is a dataset annotated several times over, so that predicted
variability can be compared with inter-reader disagreement. That disagreement,
however, is a fixed property of the data and cannot be changed.

The missing-modality setting removes this limitation. Withholding a sequence
is a deliberate change to how much evidence the model receives. Since
withholding is a deterministic transformation of the input, it cannot add
information about the annotation, so the true posterior is at least as
dispersed after removal as before. This gives a behaviour that can be
predicted in advance rather than only observed: a calibrated model should
report more uncertainty as sequences are withdrawn. Evaluating each case under
every subset makes the comparison paired, so lesion difficulty cancels and the
available evidence is the only variable that changes.

We present a three-dimensional conditional latent diffusion framework for
whole-tumour segmentation from incomplete multi-sequence MRI, and use it to
test this behaviour. Our contributions are:

\begin{enumerate}
    \item \textbf{Volumetric latent diffusion with modality dropout.} We
    extend latent diffusion segmentation to 3D multi-sequence MRI and train it
    with randomly resampled modality availability, so that one model serves
    all fifteen acquisition configurations rather than one model per protocol.

    \item \textbf{Segmentation under every acquisition configuration.} We
    report whole-tumour accuracy for all fifteen subsets of the four sequences
    on $2{,}601$ studies from BraTS 2023 and 2024, which shows how far
    performance can be maintained as sequences are withdrawn.

    \item \textbf{Uncertainty that is measured and validated under control.}
    We show that averaging predictive variance over the whole volume mixes
    confidence with lesion size, and restrict the average to the region where
    samples can disagree. Because the same case is evaluated under every
    subset, the resulting quantity can be tested against the behaviour the
    data processing inequality prescribes as evidence is removed, giving a
    paired comparison that datasets with fixed annotation ambiguity cannot
    provide.
\end{enumerate}
